% © 2026 Ing. David Jaroš — CC BY-NC-ND 4.0
%
% Rogers-Ramanujan Connection for c=3 in UBT
% Nový algebraický výsledek: θ₃(0|i)/η(i) = √2 přesně
% a jeho důsledky pro B vzorec

% UBT-AI-PROVENANCE-BEGIN
% schema: ubt-ai-provenance/v1
% tier: C_working
% ai_assistance: disclosed
% human_review: risk-based
% editorial_responsibility: Ing. David Jaroš
% policy: ../../AI_PROVENANCE.md
% notice: Working material; exhaustive human review is not claimed.
% UBT-AI-PROVENANCE-END

\documentclass[12pt,a4paper]{article}
\usepackage{amsmath,amssymb,amsthm,booktabs,tcolorbox}
\tcbuselibrary{skins,breakable}
\newtheorem{theorem}{Theorem}[section]
\newtheorem{lemma}[theorem]{Lemma}
\newtheorem{corollary}[theorem]{Corollary}
\theoremstyle{remark}
\newtheorem{remark}[theorem]{Remark}
\newcommand{\PROVED}{\textbf{[Proved]}}
\newcommand{\MC}{\textbf{[MC]}}
\newcommand{\OPEN}{\textbf{[Open]}}
\newcommand{\Lone}{\textbf{[L1]}}
\newcommand{\Lzero}{\textbf{[L0]}}
\newcommand{\STD}{\textbf{[STD]}}
\newcommand{\OBS}{\textbf{[Observation]}}

\title{\textbf{Rogers--Ramanujan Structure for $c=3$ and the UBT $B$-Formula}\\
{\large An Algebraic Derivation via $\vartheta_3(0|i)/\eta(i) = \sqrt{2}$}}
\author{Ing.\ David Jaro\v{s}}
\date{2026-05-11}

\begin{document}
\maketitle

\begin{abstract}
We establish the exact algebraic identity
$\vartheta_3(0|i)/\eta(i) = \sqrt{2}$
using Ramanujan's formula $\vartheta_3(0|i) = \pi^{1/4}/\Gamma(3/4)$,
the standard $\eta(i) = \Gamma(1/4)/(2\pi^{3/4})$,
and the reflection formula $\Gamma(1/4)\Gamma(3/4) = \pi\sqrt{2}$.
This gives a purely algebraic rewriting of the UBT candidate formula
$B = N_{\mathrm{eff}}^{3/2}(2\eta(i))^{1/4}$ as
\[
B = 12^{3/2}\cdot 2^{1/8}\cdot\left(\frac{\pi^{1/4}}{\Gamma(3/4)}\right)^{1/4},
\]
connecting it directly to Ramanujan's CM value of $\vartheta_3$ at $\tau=i$
and to the Rogers--Ramanujan--type structure of the $c=3$ partition function
$Z_{c=3} = [\vartheta_3/\eta]^3 = 2\sqrt{2}$.
The formula remains a \OBS\ (numerical accuracy $0.0066\%$);
the missing theorem is identified precisely.
\end{abstract}

%──────────────────────────────────────────────────────────────────────────────
\section{The Core Algebraic Identity \STD}
\label{sec:core}
%──────────────────────────────────────────────────────────────────────────────

\begin{theorem}[Ramanujan CM values at $\tau=i$]
\label{thm:ramanujan}
The following exact algebraic relations hold:
\begin{align}
  \vartheta_3(0|i) &= \frac{\pi^{1/4}}{\Gamma(3/4)},
  \label{eq:theta3i}\\
  \eta(i) &= \frac{\Gamma(1/4)}{2\pi^{3/4}}.
  \label{eq:etai}
\end{align}
\end{theorem}

\begin{proof}
Both are classical results of Ramanujan/Weber, proved via the
theory of elliptic integrals at the CM point $\tau=i$
(see Berndt, \emph{Ramanujan's Notebooks}, Part~III, Chapter~17).
\end{proof}

\begin{lemma}[Exact ratio \STD]
\label{lem:ratio}
\begin{equation}
  \frac{\vartheta_3(0|i)}{\eta(i)} = \sqrt{2}.
  \label{eq:ratio}
\end{equation}
\end{lemma}

\begin{proof}
From Theorem~\ref{thm:ramanujan}:
\[
\frac{\vartheta_3(0|i)}{\eta(i)}
= \frac{\pi^{1/4}/\Gamma(3/4)}{\Gamma(1/4)/(2\pi^{3/4})}
= \frac{2\pi}{\Gamma(1/4)\,\Gamma(3/4)}.
\]
The reflection formula $\Gamma(z)\Gamma(1-z)=\pi/\sin(\pi z)$ at $z=1/4$
gives
\[
\Gamma(1/4)\,\Gamma(3/4) = \frac{\pi}{\sin(\pi/4)} = \pi\sqrt{2}.
\]
Therefore $\vartheta_3(0|i)/\eta(i) = 2\pi/(\pi\sqrt{2}) = \sqrt{2}$.
\end{proof}

\begin{remark}[Numerical verification]
Direct computation confirms $\vartheta_3(0|i)/\eta(i) = 1.41421356\ldots = \sqrt{2}$
to $10^{-10}$ relative precision.
\end{remark}

%──────────────────────────────────────────────────────────────────────────────
\section{The $c=3$ Partition Function \STD}
\label{sec:c3}
%──────────────────────────────────────────────────────────────────────────────

For a CFT of three free compact bosons (central charge $c=3$),
the partition function on the self-dual torus $\tau=i$ is
\begin{equation}
  Z_{c=3}(\tau=i)
  = \left[\frac{\vartheta_3(0|i)}{\eta(i)}\right]^3
  = (\sqrt{2})^3
  = 2\sqrt{2}.
  \label{eq:Zc3}
\end{equation}

\begin{remark}[Rogers--Ramanujan context]
For $c=1/5$ (Yang--Lee model), the Rogers--Ramanujan identities give the
partition function as a ratio of $q$-products at specific modular levels.
For $c=1$ (free boson at self-dual radius), $\vartheta_3/\eta = \sqrt{2}$
is the single-boson partition function (Lemma~\ref{lem:ratio}).
For $c=3$ (three free bosons), $Z_{c=3} = (\sqrt{2})^3 = 2\sqrt{2}$.
The structure is thus a direct $c$-fold power of the $c=1$ case,
analogous to how Rogers--Ramanujan identities organise characters
by central charge, but at the distinguished CM point $\tau=i$.
\end{remark}

%──────────────────────────────────────────────────────────────────────────────
\section{Rewriting the $B$-Formula \OBS}
\label{sec:B}
%──────────────────────────────────────────────────────────────────────────────

The UBT candidate formula (numerical observation, accuracy $0.0066\%$):
\begin{equation}
  B = N_{\mathrm{eff}}^{3/2}\,(2\eta(i))^{1/4}
  = 12^{3/2}\,(2\eta(i))^{1/4}.
  \label{eq:Bcandidate}
\end{equation}

Using $\eta(i) = \vartheta_3(0|i)/\sqrt{2}$ from Lemma~\ref{lem:ratio}:
\begin{equation}
  2\eta(i) = \sqrt{2}\,\vartheta_3(0|i),
\end{equation}
so
\begin{equation}
  (2\eta(i))^{1/4} = 2^{1/8}\,\vartheta_3(0|i)^{1/4}.
\end{equation}

Substituting into~\eqref{eq:Bcandidate}:
\begin{tcolorbox}[colback=blue!3!white, colframe=blue!55!black,
  title={Algebraic rewriting of the $B$-formula}]
\begin{equation}
  \boxed{
  B = 12^{3/2}\cdot 2^{1/8}\cdot\vartheta_3(0|i)^{1/4}
    = 12^{3/2}\cdot 2^{1/8}\cdot
      \left(\frac{\pi^{1/4}}{\Gamma(3/4)}\right)^{1/4}.
  }
  \label{eq:Btheta3}
\end{equation}
\end{tcolorbox}

\paragraph{Numerical check.}
\[
12^{3/2} = 41.5692,
\quad
2^{1/8} = 1.09051,
\quad
\vartheta_3(0|i)^{1/4} = 1.08643^{1/4} = 1.02093,
\]
\[
B = 41.5692 \times 1.09051 \times 1.02093 = 46.2809.
\]
Required value: $B_{\mathrm{req}} = 2\times137/(\ln 137+1) = 46.2839$.
Relative error: $0.0066\%$. $\checkmark$

%──────────────────────────────────────────────────────────────────────────────
\section{Physical Interpretation and the Missing Theorem}
\label{sec:missing}
%──────────────────────────────────────────────────────────────────────────────

\subsection{What the formula says physically}

Equation~\eqref{eq:Btheta3} says that $B$ is determined by:
\begin{enumerate}
  \item $12^{3/2} = N_{\mathrm{eff}}^{3/2}$: the three-halves power of the
        effective mode count (algebraic structure of $\mathbb{C}\otimes\mathbb{H}$).
  \item $2^{1/8}$: a small numerical factor from the normalisation
        $\vartheta_3/\eta = \sqrt{2}$ at the self-dual point.
  \item $\vartheta_3(0|i)^{1/4}$: the quarter-power of Ramanujan's
        CM value of the Jacobi theta function at $\tau=i$.
\end{enumerate}

The factor $\vartheta_3(0|i) = \pi^{1/4}/\Gamma(3/4)$ is the
\emph{partition function of a single winding mode at the self-dual point},
and its appearance in $B$ means that $B$ is sensitive to the
self-dual geometry of the $S^1_\psi$ compactification.

\subsection{Connection to UBT winding partition function}

From \texttt{research\_tracks/quantum\_ubt/ncg\_poisson\_b0derivation.tex}:
\[
Z_{\mathrm{wind}}(\tau) = \vartheta_3(0|i\tau).
\]
At the self-dual point $\tau=1$ (i.e.\ $Z_{\mathrm{wind}}(1) = \vartheta_3(0|i)$),
the winding partition function equals Ramanujan's CM value.
Equation~\eqref{eq:Btheta3} then reads:
\begin{equation}
  B = N_{\mathrm{eff}}^{3/2}\cdot 2^{1/8}\cdot Z_{\mathrm{wind}}(1)^{1/4}.
  \label{eq:Bzwind}
\end{equation}

\subsection{The minimal missing theorem}

\begin{tcolorbox}[colback=red!4!white, colframe=red!60!black,
  title={Minimal Missing Theorem — Gap G137-B (sharpened)}]
Prove from $S[\Theta]$ that the effective coupling $B$ in
$V_{\mathrm{eff}}(n) = n^2 - Bn\ln n$ equals
\[
B = N_{\mathrm{eff}}^{3/2}\cdot 2^{1/8}\cdot Z_{\mathrm{wind}}(1)^{1/4},
\]
where $Z_{\mathrm{wind}}(1) = \vartheta_3(0|i)$ is the winding partition
function at the self-dual radius, $N_{\mathrm{eff}}$ is the effective
mode count from the algebraic structure of $\mathbb{C}\otimes\mathbb{H}$,
and all three factors arise without using $\alpha_{\mathrm{exp}}$ or
$B_{\mathrm{phenom}}$ as inputs.
\end{tcolorbox}

The sharpened form makes explicit that the three factors have
different physical origins and must be derived separately:
$N_{\mathrm{eff}}^{3/2}$ from the gauge-group structure (already [L0]),
$2^{1/8}$ from the self-dual normalisation $\vartheta_3/\eta=\sqrt{2}$ (now [STD]),
and $Z_{\mathrm{wind}}(1)^{1/4}$ from the winding spectrum at self-dual radius
(connection to $V_{\mathrm{eff}}$ is the core open problem).

%──────────────────────────────────────────────────────────────────────────────
\section{Summary}
\label{sec:summary}
%──────────────────────────────────────────────────────────────────────────────

\begin{tcolorbox}[colback=green!4!white, colframe=green!55!black,
  title={Results summary}]
\begin{center}
\begin{tabular}{lll}
\toprule
\textbf{Result} & \textbf{Level} & \textbf{Note} \\
\midrule
$\vartheta_3(0|i) = \pi^{1/4}/\Gamma(3/4)$ & \STD & Ramanujan \\
$\eta(i) = \Gamma(1/4)/(2\pi^{3/4})$ & \STD & classical \\
$\vartheta_3(0|i)/\eta(i) = \sqrt{2}$ & \Lzero/\STD & algebraic identity \\
$Z_{c=3}(i) = 2\sqrt{2}$ & \STD\ (given $c=3$) & partition function \\
$B = 12^{3/2}\cdot 2^{1/8}\cdot\vartheta_3(0|i)^{1/4}$ & \OBS & $0.0066\%$ accuracy \\
Physical derivation of $B$ from $S[\Theta]$ & \OPEN & Gap G137-B \\
\bottomrule
\end{tabular}
\end{center}
\end{tcolorbox}

\end{document}
